\section{主观题}
\begin{enumerate}
  \item 不会考的太前沿、深入、刁钻、细化。基础理论的掌握。
  \item 民商诉连接点。
        \begin{enumerate}
          \item 商法 = 公司法。出资责任（股转）、决议效力、董监高信义义务。
          \item 民法。
                \begin{enumerate}
                  \item 合同着手。排除人格、家庭、婚姻、继承、收养。
                  \item 合同之债、物权变动（合同 + 登记）、侵权之债（归责原则）。
                  \item 民事法律行为、代理 \yjt 解决合同效力问题。
                  \item 抵押、保证。
                  \item 买卖合同、租赁合同、建设工程合同。
                  \item 合同履行、保全（代位之诉、撤销之诉）、转让、解除、责任（违约、缔约过失）。
                \end{enumerate}
          \item 民诉。当事人制度、诉讼/仲裁选择、管辖、诉的种类。
        \end{enumerate}
  \item 请求权分析法。
        \begin{enumerate}
          \item 债权。
                \begin{enumerate}
                  \item 合同。缔约过失（信赖利益损失。订约费用、机会）、违约（继续履行、补救措施、损害赔偿金、违约金、定金）。
                  \item 无因管理/准合同。必要费用、所负债务、所受损害。
                  \item 不当得利。返还问题。基于给付、非基于给付。
                  \item 侵权。赔偿问题。以违法为前提。停止侵害、排除妨害、消除危险、消除影响、恢复名誉、赔礼道歉。
                  \item 单方允诺。
                \end{enumerate}
          \item 物权。返还原物、排除妨害、消除危险、恢复原状。
          \item 占有。占有返还、停止侵害、排除妨害、消除危险、必要费用。
          \item 合同 \dy 无因管理 \dy 物权/占有 \dy 不当得利 \dy 侵权。
          \item 每个请求权都是分析 3 个问题。
                \begin{enumerate}
                  \item 内容 \yjt 请求权是什么。
                  \item 成立 \yjt 构成要件。
                  \item 不成立 \yjt 抗辩事由。
                \end{enumerate}
        \end{enumerate}
  \item 答题。
        \begin{enumerate}
          \item 答背后的原理即可。关键词、范式（什么对什么有权利，因为什么）。不用先说原理，再结合本案怎样怎样。\al{生效。所附条件已成就}。
          \item 不要超过一行。\al{占有人对侵占人享有占有返还请求权}。
          \item \al{不能。因为流押条款无效}。
          \item 行为、主体、权利的表述法言法语。
          \item 卷面格局。
          \item 题纲目录。开题（背景、本质、意义）、分析、结论。
        \end{enumerate}
  \item 怎么做题。近 10 年真题。
  \item \bs{问你能不能 xx、能否、是否属于，答全要件}。
\end{enumerate}
